% Source-preserving integrated proof body; input from toe_round7_proofs.tex.
\subsubsection*{Round 7: a Lorentz-covariant free curvature field on the Round 6 Hilbert space}
\label{sec:toe-r7-covariant}
\noindent\emph{The displayed arguments below are the proof-bearing source.  The current
repository regression is \veri{v1032\_\allowbreak covariant\_\allowbreak curvature.py}; its finite checks audit the
implementation and do not replace the universal steps in the proof.}

\begingroup
\def\TT{\mathrm{TT}}
\def\cH{\mathcal H}
\begin{keybox}[Scope and exact claim boundary]
The Round 6 target gave positive quantum electric-curvature variables
$E,F$ and a microcausal continuum limit, but did not construct the full
four-dimensional Riemann field or its Lorentz boosts. We close that specific
free-field gap. A polynomial covariant curvature kernel equals a positive
two-polarization kernel on the null cone. It defines a positive Gaussian
Wightman field, the vacuum linearized Riemann (equivalently Weyl) tensor,
with a strongly continuous positive-energy Poincare representation and
helicities $+2,-2$. Its electric part is exactly $E/2$ and its time derivative
is $F/2$; it adds no one-particle states. The construction is a free continuum
extension of the target, not the missing derivation from the TFPT parent.
It proves neither nonlinear gravity, local-net equality with the smaller
$E,F$ system, a local lattice realization of every curvature component,
nor a complete T7 or TOE. No RH statement is involved.
\end{keybox}

\paragraph{Conventions and the algebraic bridge.}
Use $\eta=\operatorname{diag}(-1,1,1,1)$ and a covariant null momentum
\[
 k_\mu=(-r,p_1,p_2,p_3),\qquad r=|p|>0,
 \qquad e^{ik_\mu x^\mu}=e^{-irt+ip\cdot\mathbf x}.
\]
All tensor components in this note carry lower indices. For a symmetric
metric amplitude $h$, define the linear curvature map
\begin{equation}
 (A(k)h)_{\mu\nu\rho\sigma}
 =\frac12\left(k_\sigma k_\nu h_{\mu\rho}
 +k_\rho k_\mu h_{\nu\sigma}
 -k_\rho k_\nu h_{\mu\sigma}
 -k_\sigma k_\mu h_{\nu\rho}\right).
 \label{eq:toe-r7-covariant-A}
\end{equation}
It is the Fourier transform of
\[
 R_{\mu\nu\rho\sigma}
 =\tfrac12(\partial_\rho\partial_\nu h_{\mu\sigma}
 +\partial_\sigma\partial_\mu h_{\nu\rho}
 -\partial_\sigma\partial_\nu h_{\mu\rho}
 -\partial_\rho\partial_\mu h_{\nu\sigma}).
\]
No positive covariant quantum metric potential is assumed. Instead introduce
the numerical tensor
\begin{equation}
 \Pi_{\mu\nu,\alpha\beta}
 =\tfrac12(\eta_{\mu\alpha}\eta_{\nu\beta}
 +\eta_{\mu\beta}\eta_{\nu\alpha}
 -\eta_{\mu\nu}\eta_{\alpha\beta}),
 \qquad K(k)=A(k)\Pi A(k)^T.
 \label{eq:toe-r7-covariant-K}
\end{equation}
In a Euclidean orthonormal basis of symmetric $4\times4$ matrices, $\Pi$
has six positive and four negative eigenvalues. It is not a physical
Hilbert-space covariance by itself.

Let $S$ embed a spatial symmetric matrix in the $ij$ block of a spacetime
matrix, with $h_{0\mu}=0$. The orthogonal spatial transverse-traceless
projector is
\begin{align}
 P_{ij}&=\delta_{ij}-p_ip_j/r^2,\\
 (P_\TT h)_{ij}&=P_{ia}P_{jb}h_{ab}
       -\tfrac12 P_{ij}P_{ab}h_{ab}.
 \label{eq:toe-r7-covariant-PTT}
\end{align}
It is an orthogonal rank-two projector on $\operatorname{Sym}^2(\mathbb R^3)$.

\begin{lemma}[Null-cone positive factorization]
For every nonzero null $k=(-r,p)$,
\begin{equation}
 K(k)=A(k)S P_\TT(p)S^TA(k)^T
       =L(k)L(k)^T,\qquad L(k)=A(k)S P_\TT(p).
 \label{eq:toe-r7-covariant-positive}
\end{equation}
Consequently $K(k)$ is positive semidefinite of rank exactly two.
It is nevertheless a polynomial of degree four in all four components of $k$.
\end{lemma}
\begin{proof}
Substitution in~\eqref{eq:toe-r7-covariant-A} gives $A(k)(k\xi^T+\xi k^T)=0$ for every
$\xi$. The identity~\eqref{eq:toe-r7-covariant-positive} can be checked with $p=(0,0,r)$:
the two surviving transverse orthonormal matrices are
\[
 h_+=\operatorname{diag}(1,-1,0)/\sqrt2,\qquad
 h_\times=(e_1e_2^T+e_2e_1^T)/\sqrt2,
\]
and direct contraction gives
$A\Pi A^T=(ASh_+)(ASh_+)^T+(ASh_\times)(ASh_\times)^T$.
Spatial rotational covariance and degree-four homogeneity give the result
for every nonzero $p$. The map is injective on the two-dimensional TT
space because $(A(k)Sh)_{0i0j}=r^2h_{ij}/2$. This proves rank two.
Polynomiality follows directly from~\eqref{eq:toe-r7-covariant-K}; no division by $r$
occurs there.

For an independent algebraic certificate, the checker uses ten orthonormal
metric components and 21 pair-symmetric bivector components. Write
$\omega$ as an independent variable and $C(p)=r^4P_\TT(p)$, the polynomial
Round 6 bracket. It divides each of the 441 entries of
\[
 A(-\omega,p)\bigl(r^4\Pi-SC(p)S^T\bigr)A(-\omega,p)^T
\]
by $\omega^2-r^2$ and verifies an identically zero polynomial remainder.
The off-shell identity and the identity without the trace-reversal term
are separately required to fail. These are algebraic checks, not sampling
of momentum directions or a numerical positivity test.
\end{proof}

The same map obeys both pair antisymmetries, pair interchange, the algebraic
Bianchi identity $R_{\mu[\nu\rho\sigma]}=0$, and the differential Bianchi
identity $k_{[\tau}R_{|\mu\nu|\rho\sigma]}=0$.
On the TT null range the Ricci contraction vanishes,
$\eta^{\mu\rho}R_{\mu\nu\rho\sigma}=0$. For example, insert $h_{0\mu}=0$,
$p_ih_{ij}=0$, $h_{ii}=0$, and $k^2=0$ into~\eqref{eq:toe-r7-covariant-A}; every contraction
contains one of those four zero factors. Thus the constructed vacuum
curvature is also its linearized Weyl tensor. These identities hold as
operator-valued distributional identities after the null quotient below.

\paragraph{Positive field, local commutator, and boosts.}
Index the independent pair-symmetric bivector entries by $I,J$.
Define the two-point distribution
\begin{equation}
 W_{IJ}(x-y)=\int_{\mathbb R^3}
 \frac{d^3p}{(2\pi)^3\,2r}\,
 e^{-ir(t_x-t_y)+ip\cdot(\mathbf x-\mathbf y)}K_{IJ}(-r,p).
 \label{eq:toe-r7-covariant-W}
\end{equation}
The Bianchi identities make one of the 21 algebraic components redundant;
retaining it with its null relation avoids an arbitrary choice of twenty
coordinates.

\begin{theorem}[Free covariant curvature field]
Equation~\eqref{eq:toe-r7-covariant-W}, with Wick's rule for higher vacuum correlations,
defines a positive free bosonic operator-valued tempered distribution on
a symmetric Fock space. Its commutators vanish at spacelike separation.
There is a strongly continuous unitary representation of the proper
orthochronous Poincare group, with invariant cyclic Fock vacuum, spectrum
in the closed forward cone, and covariant rank-four curvature transformation.
Its only one-particle helicities are $+2$ and $-2$.
\end{theorem}
\begin{proof}
For Schwartz test tensors $f$, factorization~\eqref{eq:toe-r7-covariant-positive} gives
\begin{equation}
 \|f\|_1^2=\int\frac{d^3p}{(2\pi)^3\,2r}
 \widehat f(k)^*K(k)\widehat f(k)
 =\int\frac{d^3p}{(2\pi)^3\,2r}
 \|L(k)^T\widehat f(k)\|^2\geq0.
 \label{eq:toe-r7-covariant-norm}
\end{equation}
Quotient the zero-norm tensors and complete to $\cH_1$.
Write $u_f$ for the one-particle image of the test tensor. With $a(u)$
antilinear in $u$, set $R(f)=a^\dagger(u_f)+a(u_{\bar f})$ on the
finite-particle domain of $\Gamma_s(\cH_1)$, using the Fourier convention
implicit in~\eqref{eq:toe-r7-covariant-W}. This is complex linear in $f$ and obeys
$R(f)^*=R(\bar f)$ on that domain. It gives~\eqref{eq:toe-r7-covariant-W}. The usual creation
estimate $\|a^\dagger(u)\psi_n\|\leq\sqrt{n+1}\|u\|\|\psi_n\|$
and continuity of~\eqref{eq:toe-r7-covariant-norm} make this an operator-valued tempered
distribution. Wick positivity follows from this explicit Fock
construction, not from assuming $\Pi$ is positive.

Infrared convergence and temperedness are immediate: $K=O(r^4)$,
so $K/(2r)=O(r^3)$ at zero and has only polynomial growth at infinity.
Together with the $d^3p$ measure the radial integrand is $O(r^5)$ near zero.
Schwartz test functions control every ultraviolet moment.

Let
\[
 D_0(x)=\int\frac{d^3p}{(2\pi)^3\,2r}
 (e^{-irt+ip\cdot\mathbf x}-e^{irt-ip\cdot\mathbf x})
\]
be the scalar massless commutator distribution. Since $K(-k)=K(k)$,
the curvature commutator is the constant-coefficient differential operator
$K_{IJ}(-i\partial)D_0(x-y)$. The scalar commutator has light-cone support,
and differentiation does not enlarge support. This proves locality of the
curvature fields, despite the nonlocal spatial TT projector used to exhibit
positivity.

All indices in~\eqref{eq:toe-r7-covariant-A}--\eqref{eq:toe-r7-covariant-K} are tensorial and $\Pi$ is built
only from $\eta$. Therefore
\begin{equation}
 K(\Lambda k)=D^{(4)}(\Lambda)K(k)D^{(4)}(\Lambda)^T
 \label{eq:toe-r7-covariant-boost}
\end{equation}
for the covector action of every proper orthochronous Lorentz matrix.
The positive-null-cone measure is invariant. The induced contragredient
test-tensor transformation, together with translations, consequently
preserves~\eqref{eq:toe-r7-covariant-norm} and its null space. It gives a unitary
one-particle representation. The test-tensor action is continuous in the
Schwartz topology, so the norm estimate gives strong continuity on a
dense set and then on the completion. Second quantization gives the claimed
Fock representation. Translations act by positive-energy null momenta;
finite sums and closure keep the many-particle spectrum in the forward cone.
The null cone tip has zero one-particle measure. Thus every nonvacuum
$n$-particle wave function has strictly positive total energy almost
everywhere, and the only zero-energy vector in this constructed Fock
representation is the vacuum. This is not a theorem of uniqueness across
all representations or all possible TFPT states.

For helicity, at $k=(-1,0,0,1)$ rotations about the third axis act on
$(h_+,h_\times)$ by
\[
 \begin{pmatrix}\cos2\theta&-\sin2\theta\\
                 \sin2\theta&\cos2\theta\end{pmatrix}.
\]
Hence $h_+\pm i h_\times$ have phases $e^{\mp2i\theta}$.
The translation part of the null little group is trivial on the curvature:
put $q=(u^2+v^2)/2$ and
\[
 \Lambda(u,v)=
 \begin{pmatrix}
 1+q&-u&-v&q\\
 -u&1&0&-u\\
 -v&0&1&-v\\
 -q&u&v&1-q
 \end{pmatrix}.
\]
It preserves $\eta$ and $k$, and sends the transverse covectors
$e_1\mapsto e_1+u k$, $e_2\mapsto e_2+v k$. Each polarization change is
therefore $k\xi^T+\xi k^T$, which $A(k)$ annihilates.
This proves the two helicities without assuming away the noncompact
little-group translations. The checker verifies these identities for
symbolic $u,v$ and all rotation angles through a rational parametrization.
\end{proof}

\paragraph{It is the same free target, with a larger covariant field family.}
The Round 6 continuum construction has, on the physical TT range,
\[
 E=r^2q_\TT,\qquad F=r^2p_\TT,\qquad
 \dot q_\TT=p_\TT,\qquad\dot p_\TT=-r^2q_\TT,
\]
with $C=r^4P_\TT$. In radiation gauge~\eqref{eq:toe-r7-covariant-A} gives exactly
\begin{equation}
 R_{0i0j}=\frac{r^2}{2}q_{ij}=\frac12E_{ij},
 \qquad\partial_tR_{0i0j}=\frac12F_{ij}.
 \label{eq:toe-r7-covariant-same}
\end{equation}
Consequently the spectral electric-electric covariance is
$C/(8r)$, electric-time-derivative covariance is $iC/8$ at equal time
in the order $\langle R\dot R\rangle$, and derivative-derivative covariance
is $rC/8$. These are precisely one quarter of the Round 6 $EE,EF,FF$
covariances, respectively, including the nonsymmetrized commutator part.

For completeness this matching is more than equality of a few entries.
Both one-particle spaces identify with the completion of spatial TT
amplitudes in $L^2(d^3p/((2\pi)^3 2r))$. The map from curvature tests is
$f\mapsto P_\TT S^TA^T\widehat f$. Electric tests alone have dense image:
on every annulus $0<\epsilon<r<R$ their map is multiplication by the
invertible scalar $r^2/2$ followed by $P_\TT$. Every smooth compactly
supported TT amplitude on such an annulus can be extended to a spacetime
Schwartz test in a tubular neighbourhood of the positive null cone.
Cutting off at zero and infinity is dense in the one-particle norm.
Thus~\eqref{eq:toe-r7-covariant-same} identifies the completed one-particle spaces, and
the associated Fock spaces, with no extra physical oscillator or spectator.

There is an important locality qualification. At a fixed time,
\[
 R_{0ijk}=\tfrac12(\partial_k p_{ij}-\partial_j p_{ik})
          =\frac1{2r^2}(\partial_kF_{ij}-\partial_jF_{ik})
\]
in Fourier notation; the spatial components similarly contain two spatial
derivatives of $q=E/r^2$. Inverse powers of $r^2$ cancel in the covariant
commutator kernel, but are still present in this reconstruction from $E,F$.
The global Hilbert-space identification does not by itself prove equality
of local observable nets, nor that every $R$ component has a bounded-range
lattice approximant in the previous $E,F$ algebra. No such claim is used
in the positivity, covariance, or locality theorem above.

\paragraph{Validation, provenance, and the remaining obligation.}
The repository certificate \veri{v1032\_\allowbreak covariant\_\allowbreak curvature.py} is standalone
SymPy code. Its 34 exact checks include the all-momentum polynomial identity,
gauge and Bianchi identities, null Ricci contraction, positivity
factorizations, a nontrivial rational Lorentz boost, the full null little
group, and the exact electric and mixed time-derivative covariance normalization. Four rational
fibres supplement, but do not replace, the universal proof. No floating-point
eigenvalue scan is used. The distributional/Fock completion rests on the
written norm, continuity, density and support arguments, not on a finite
matrix test pretending to construct an infinite-dimensional field.

The inputs are the declared Round 6 free target and ordinary linearized
curvature algebra. No microscopic TFPT Hamiltonian is imported or derived
by this checker; only the repository's reporting harness is shared.
In particular this note does not manufacture the required
shared parent: it does not prove that the microscopic seam, charged matter,
mirror sector and gravity jointly flow to this field or to its nonlinear
completion. A finite-spacing Lorentz-breaking regulator is not asserted
to be exactly boost invariant. The vacuum homogeneous-block choice and
the interacting constraint/renormalization problems remain separate.

For historical context, massless-spin field representations and the role
of the non-semisimple little group are discussed in Weinberg's primary
paper\footnote{S. Weinberg, ``Feynman Rules for Any Spin. II. Massless
Particles'', Phys. Rev. 134, B882 (1964),
\url{https://doi.org/10.1103/PhysRev.134.B882}. The accessible abstract was
used for context; no inaccessible theorem is imported as a proof.}.
The result here is an explicit match of that free-field structure to the
existing target, not a claim that free helicity-two fields are a new
physical theory. T1--T8 retain their existing open status.
\endgroup
